\documentclass[11pt,a4paper]{article}

% Lingua e caratteri (compilare con XeLaTeX)
\usepackage{fontspec}
\usepackage[italian]{babel}

\usepackage[margin=2.5cm]{geometry}
\usepackage{amsmath, amssymb}
\usepackage{graphicx}
\usepackage{enumitem}
\usepackage[hidelinks]{hyperref}

% Dati del compito
\newcommand{\corso}{Nome del corso}
\newcommand{\compito}{Compito 1}
\newcommand{\studente}{Nome Cognome}
\newcommand{\docente}{Nome del docente}

% Domanda numerata e ambiente per la risposta
\newcounter{domanda}
\newcommand{\domanda}[1]{%
  \stepcounter{domanda}%
  \section*{Domanda \thedomanda\ifx\relax#1\relax\else: #1\fi}}
\newenvironment{risposta}{\par\noindent\textbf{Risposta.}\ }{\hfill$\blacksquare$\par}

\begin{document}

\noindent
\begin{minipage}[t]{0.6\textwidth}
  \textbf{\corso}\\
  \compito
\end{minipage}%
\begin{minipage}[t]{0.4\textwidth}
  \raggedleft
  \studente\\
  Docente: \docente\\
  \today
\end{minipage}
\par\medskip\hrule\medskip

\domanda{Derivate}
Calcolare $\dfrac{d}{dx}\left(x^2 \sin x\right)$.

\begin{risposta}
  Per la regola del prodotto,
  \[
    \frac{d}{dx}\left(x^2 \sin x\right) = 2x \sin x + x^2 \cos x.
  \]
\end{risposta}

\domanda{}
Enunciare due proprietà dell’integrale definito.

\begin{risposta}
  \begin{enumerate}[label=(\alph*)]
    \item Linearità: $\int_a^b (\alpha f + \beta g) = \alpha \int_a^b f + \beta \int_a^b g$.
    \item Additività: $\int_a^b f = \int_a^c f + \int_c^b f$.
  \end{enumerate}
\end{risposta}

\end{document}
